\documentclass[10pt,a4paper]{amsart}
\usepackage[margin=19mm]{geometry}
\usepackage{amsmath,amssymb,mathtools}
\usepackage[scr=rsfs,cal=euler]{mathalfa}
\usepackage{xfrac}
\usepackage[unicode,hidelinks,bookmarks=false]{hyperref}
\newcommand{\ie}{i.e.\ }
\newcommand{\cf}{cf.\ }
\newcommand{\resp}{respectively\ }
\newcommand{\Subsection}{\subsection}
\providecommand{\nobreakdash}{\nobreak\hbox{-}}
\setlength{\emergencystretch}{2em}
\multlinegap0pt
\usepackage{amssymb}
\usepackage{mathtools}
\usepackage[normalem]{ulem}
\usepackage{xcolor}
\usepackage[shortlabels]{enumitem}
\usepackage{tcolorbox}
\newtheorem{proposition}{Proposition}
\newcommand{\G}{{\mathbb G}}
\newcommand{\PP}{{\mathbb P}}
\newcommand{\ko}{{\mathcal O}}
\title{\bf Nash Loci}
\author{Luca Sodomaco, Julian Weigert}
\date{Source: arXiv:2608.27300v1 (CC BY 4.0)}
\begin{document}
\maketitle

\noindent\emph{Source and licence.} \bf Nash Loci. Version arXiv:2608.27300v1 (CC BY 4.0), distributed by arXiv under the Creative Commons Attribution 4.0 International licence, http://creativecommons.org/licenses/by/4.0/, as recorded in the arXiv licence metadata for this exact version and shown on the abstract page https://arxiv.org/abs/2608.27300v1. Full licence notice: \texttt{licenses/arxiv-2608.27300-CC-BY-4.0.txt}.\par\smallskip

\noindent\textit{Editorial scope.} Counted unit: the article's proposition variety (source0.tex lines 879--888). The article's complete original proof of it is delivered with it (source0.tex lines 889--898, one \texttt{proof} environment of 10 lines). No mathematics is written, repaired or re-derived: the statement and the proof are the article's own lines, in the article's own notation, and no retained line is substituted or re-flowed.  Retained uncounted as the source-local context the proof needs: lines 726--732; lines 813--815; lines 824--826.  Nothing was omitted from any retained span, so the package carries no omission record.  No retained line contains a citation, so the wrapper carries no bibliography and no bibliography entry.  No retained line contains a figure environment, an image-inclusion command or drawing code, so no image is delivered and the package declares no asset dependency.  Editorial formatting only, in addition: an AMS \texttt{amsart} preamble that restates the article's own declarations and macros used by the retained lines, namely the environments \texttt{proposition} on separate counters, together with the standard packages those lines need.  The retained environments print in this document as Proposition 1 in order of appearance, whereas the article prints the same items with the numbering of its own full document.  The complete native source of the article as distributed in the arXiv e-print for this exact version is bundled unchanged as \texttt{sources/208706/source0.tex}.
\begin{proposition}\label{prop: when Chow is of expected codimension}
Under the previous notations, we have that $\mathrm{Ch}_\alpha(X)$ has the expected codimension $e$ if and only if for every nonempty $I \subseteq [n]$ we have
\begin{equation}\label{eq:genericity_conditions}
    \dim p_I(X) \ge |\beta_I|-e
\end{equation}
where $p_I(X)$ denotes the projection of $X$ onto $\prod_{i\in I}\mathbb{P}^{k_i}$.
\end{proposition}

\begin{equation}\label{eq: ith tautological sequence}
    0\to S_i \to \ko_{\G_i}^{\oplus (k_i+1)} \to Q_i \to 0\,,
\end{equation}

\begin{equation}\label{eq: multidegree X}
    [X]\ =\ \sum_{|\gamma|=c} \delta_\gamma(X)\, H_1^{\gamma_1}\cdots H_n^{\gamma_n}\,.
\end{equation}

\begin{proposition}\label{prop: class of multigraded associated variety}
Assume that the inequalities in Proposition~\ref{prop: when Chow is of expected codimension} are satisfied. Then
\[
    [\mathrm{Ch}_\alpha(X)] \ =\  \pi_*\,p^*[X]\in A^*[\G_\alpha]\,.
\]
Equivalently, using the multidegree expansion \eqref{eq: multidegree X},
\begin{equation}\label{eq: deg multigraded associated variety}
    [\mathrm{Ch}_\alpha(X)]\ =\ \sum_{\substack{|\gamma|=c\\\gamma_i\ge\alpha_i\ \forall i}} \delta_\gamma(X)\ \prod_{i=1}^n s_{\gamma_i-\alpha_i}(S_i)\ =\ \sum_{\substack{|\gamma|=c\\\gamma_i\ge \alpha_i\ \forall i}}\delta_\gamma(X)\ \prod_{i=1}^n c_{\gamma_i-\alpha_i}(Q_i)\,.
\end{equation}
\end{proposition}

\begin{proof}
From the incidence definition, as a cycle class one has $[\mathrm{Ch}_\alpha(X)]=\pi_*p^*[X]$.
Using \eqref{eq: multidegree X} and $p^*(H_i)=\zeta_i$ gives
\[
    \pi_*p^*[X]\ =\ \sum_{|\gamma|=c}\delta_\gamma(X)\ \pi_*\!\left(\prod_{i=1}^n \zeta_i^{\gamma_i}\right)\,.
\]
Since $\pi$ is a fiber product of projective bundles $\PP(S_i)\to \G_\alpha$ of relative dimension $\alpha_i$, projective bundle pushforward yields $(\pi_i)_*\bigl(\zeta_i^{\alpha_i+j}\bigr)=s_j(S_i)$ ($j\ge 0$), $(\pi_i)_*(\zeta_i^m)=0$ for $m<\alpha_i$.
Thus only terms with $\gamma_i\ge \alpha_i$ contribute and $\pi_*\!\left(\prod_{i=1}^n \zeta_i^{\gamma_i}\right) = \prod_{i=1}^n s_{\gamma_i-\alpha_i}(S_i)$.
Finally $s(S_i)=c(S_i)^{-1}$ and $c(S_i)c(Q_i)=1$ from the Whitney sum formula applied to \eqref{eq: ith tautological sequence}, hence $s_{t}(S_i)=c_t(Q_i)$ for all $t$.
\end{proof}

\end{document}
